\phantomsection\label{sec:qft-renormalization-rg-dp10}

前面的微扰规则已经把散射振幅写成顶角、传播子和内部动量积分的组合。新的困难出现在圈积分的高动量端：某些积分若一直延伸到任意大的内部动量，会发散，因而“直接代入裸质量和裸耦合再积分”并不是一个完整的计算定义。解决它需要先把发散积分改写成带有临时控制参数的有限表达式，再用实验可定义的质量、耦合和场归一化重新参数化作用量，使最终可观测量不依赖这个临时参数。前一步称为正则化，后一步称为重整化。

重整化以后，同一个物理过程仍可在不同外部能标下被探测，因此有限参数必须随分辨尺度变化而保持可观测量不变；描述这种尺度依赖的微分方程就是重整化群。当实验能量远低于某些重自由度的质量尺度时，这些自由度不能成为真实外态，却仍通过虚过程留下低能效应；有效场论把这些高能效应编码成局域算符及其 Wilson 系数。

圈动量定义为顶角守恒约束后仍未固定的独立内部 Fourier 积分变量；一般内部传播线处于离壳状态。紫外发散说明“裸参数直接等于实验参数”的写法没有把短程贡献与测量条件分开。正则化引入临时数学参数；反项与重整化条件再把这种临时依赖吸收到有限个参数定义中。最终可观测量只依赖物理输入与外部运动学尺度，而不依赖正则化器本身。


共同放大所有独立圈物理四动量，\(\ell_r^\mu\mapsto\Lambda\ell_r^\mu\) 且 \(\Lambda\to\infty\)，即可诊断图的表面紫外次数。四维圈测度贡献四次幂，玻色传播子贡献负二次幂，Dirac 传播子在大动量下贡献负一次幂；若顶角分子还含 \(N_\partial\) 个圈动量因子，则连通图满足
\begin{equation}
 \boxed{
 \omega_{\rm UV}=4L-2I_B-I_F+N_\partial,}
 \label{eq:superficial-degree-general-dp10}
\end{equation}
其中 \(L\) 是独立圈数，\(I_B\) 与 \(I_F\) 分别为内部玻色子线与费米子线数。\(\omega_{\rm UV}<0\) 只说明所有圈动量共同放大时表面收敛，仍须检查发散子图；\(\omega_{\rm UV}\ge0\) 也只是最坏幂次上界，对称性还可以排除相应局域结构。

对四维 \(\varphi^4\) 理论，每个四次顶角提供四个线端。若连通图含 $V$ 个顶角、$I$ 条内部线和 $E$ 条外线，则线端计数与独立圈数满足
\begin{equation}
 \boxed{
 4V=2I+E,
 \qquad
 L=I-V+1.}
 \label{eq:phi4-topology-r68}
\end{equation}
把这两条纯拓扑关系代入\eqnrefs{eq:superficial-degree-general-dp10}，便得到
\begin{equation}
 \boxed{\omega_{\varphi^4}=4-E,}
 \label{eq:phi4-superficial-degree-dp10}
\end{equation}
所以除真空图外，只有二点与四点函数需要原理论已有的质量、场强和四次耦合局域结构作为反项。QED 的每个最小耦合顶角连接一个光子腿和两个费米子腿。相应的光子线端、费米子线端和独立圈数同时满足
\begin{equation}
 \boxed{
 V=2I_\gamma+E_\gamma,
 \qquad
 2V=2I_\psi+E_\psi,
 \qquad
 L=I_\gamma+I_\psi-V+1.}
 \label{eq:qed-topology-r68}
\end{equation}
把这些拓扑恒等式代入 $\omega=4L-2I_\gamma-I_\psi$，即可消去内部线数并得到
\begin{equation}
 \boxed{
 \omega_{\rm QED}=4-E_\gamma-\frac32E_\psi.}
 \label{eq:qed-superficial-degree-dp10}
\end{equation}
例如光子二点函数表面计数给二次发散，但 Ward 恒等式要求 \(q_\mu\Pi^{\mu\nu}=0\)，把允许张量限制到横向结构，实际只剩对数发散。幂次计数决定潜在紫外阶数，规范对称性再决定哪些局域项真正允许出现。

发散积分必须先被数学地定义。维数正则化取
\begin{equation*}
 d=4-2\epsilon,
 \qquad \epsilon>0,
\end{equation*}
并引入具有\emph{物理四动量量纲}的尺度 \(\mu_{\rm DR}\)。四维无量纲耦合的一般圈积分写成
\begin{equation}
 \boxed{
 \mu_{\rm DR}^{2\epsilon}
 \int\frac{\dd^d \ell}{(2\pi)^d}(\cdots),}
 \label{eq:dim-reg-measure-general-dp10}
\end{equation}
其中 $\ell^\mu$ 始终是物理四动量。若改用逆长度变量 $\kappa^\mu=\ell^\mu/\hbar$，则测度提供 $\hbar^{-d}$，两个二次传播子分母提供 $\hbar^4$，而维数正则化尺度 $\mu_\kappa^{2\epsilon}=(\mu_{\rm DR}/\hbar)^{2\epsilon}$ 再提供 $\hbar^{-2\epsilon}$。因为 $d=4-2\epsilon$，总幂次
\[
 -d+4-2\epsilon=0
\]
恰好为零，因此物理动量形式中的显式 $\hbar$ 由变量换元严格抵消。典型积分在 $\epsilon\to0$ 时产生局域极点与外部尺度对数；前者由反项吸收，后者保留为可测尺度依赖。

有外部四动量时，一圈分子通常还含 $\ell^\mu$、$\ell^\mu\ell^\nu$ 等张量。Lorentz 协变性保证它们只能展开在 $\eta^{\mu\nu}$ 与外部四动量张量基上，因此张量积分可以先约化为少数标量母积分；这一步称为 Passarino--Veltman 约化。沿用\eqnrefs{eq:reduced-loop-measure-si-r10} 的物理四动量测度，对任意具有动量量纲的质量参数 $M_i$ 定义
\begin{equation}
\boxed{
\begin{aligned}
 A_0(M^2)
 &\equiv \frac{16\pi^2}{\ii}\int_\ell^{(d)}
 \frac{1}{\ell^2-M^2+\ii0},\\
 B_0(p^2;M_1^2,M_2^2)
 &\equiv \frac{16\pi^2}{\ii}\int_\ell^{(d)}
 \frac{1}{(\ell^2-M_1^2+\ii0)[(\ell+p)^2-M_2^2+\ii0]},\\
 C_0(p_1^2,p_2^2,(p_1+p_2)^2;M_1^2,M_2^2,M_3^2)
 &\equiv \frac{16\pi^2}{\ii}\int_\ell^{(d)}
 \frac{1}{D_1D_2D_3},
\end{aligned}}
\label{eq:pv-scalar-masters-dp90}
\end{equation}
其中
\[
D_1=\ell^2-M_1^2+\ii0,\quad
D_2=(\ell+p_1)^2-M_2^2+\ii0,\quad
D_3=(\ell+p_1+p_2)^2-M_3^2+\ii0.
\]
按这一定义，$A_0$、$B_0$、$C_0$ 的动量量纲依次为 $2,0,-2$。例如秩二三点积分不需要引入新的非标量函数作为物理输入，因为必可写成
\begin{equation}
\frac{16\pi^2}{\ii}\int_\ell^{(d)}\frac{\ell^\mu\ell^\nu}{D_1D_2D_3}
=C_{00}\eta^{\mu\nu}+\sum_{i,j=1}^{2}C_{ij}p_i^\mu p_j^\nu .
\label{eq:pv-rank2-decomposition-dp90}
\end{equation}
分别用 $\eta_{\mu\nu}$、$p_{1\mu}p_{1\nu}$、$p_{1\mu}p_{2\nu}$、$p_{2\mu}p_{2\nu}$ 收缩即可得到 $C_{00},C_{ij}$ 的线性方程；再利用 $2\ell\!\cdot p_i$ 与相邻传播子分母之差，把右端化为\eqnrefs{eq:pv-scalar-masters-dp90}。因此 QED、QCD 与标准模型的一圈张量约化都归结为同一件事：先在 Lorentz 张量基上展开，再用传播子恒等式把张量积分精确消元为标量主积分。

重整化把同一裸作用量精确拆成“用有限参数写成的原型拉氏量 + 局域反项”。以实标量场为例，写成
\begin{equation*}
 \varphi_B=Z_\varphi^{1/2}\varphi_R,
 \qquad
 m_B^2=m_R^2+\delta m^2+\cdots,
 \qquad
 g_{4,B}=\mu_{\rm DR}^{2\epsilon}(g_4+\delta g_4),
\end{equation*}
即可得到
\begin{equation*}
 \Lag_B=\Lag_{\rm R}+\Lag_{\rm ct}.
\end{equation*}
反项系数由重整化条件以及裸理论对临时正则化参数的独立性确定。

\begin{note}[常用重整化方案]
壳上方案以传播子物理极点、留数和指定运动学点的电荷作为输入；最小减除只减去 \(1/\epsilon\) 极点；\(\overline{\rm MS}\) 方案同时减去 \(1/\epsilon-\gamma_E+\ln4\pi\)。有限阶计算会保留方案依赖，但同一可观测量在相同完整微扰阶上的方案差异属于更高阶截断误差。
\end{note}


\begin{exbox}[$\varphi^4$ 理论的三类重整化]
单靠“圈积分发散”还看不出应该重定义哪个参数。最清楚的办法是分别看二点函数与四点函数。

一圈二点 1PI 图只有一个四次顶角。两条外腿占用顶角的两个场，剩余两个场彼此收缩成蝌蚪圈，因此
\begin{equation}
 \boxed{
 -\ii\Sigma_{\rm tad}^{(1)}
 =\frac{-\ii g_4}{2}
 \mu_{\rm DR}^{2\epsilon}
 \int\frac{\dd^d \ell}{(2\pi)^d}
 \frac{\ii}{\ell^2-m_R^2c^2+\ii0^+}.}
 \label{eq:phi4-tadpole-selfenergy-dp52}
\end{equation}
$1/2$ 是同一顶角上两条内部腿互换的对称性因子。这个积分不含外部四动量 $p$，所以它只能平移传播子极点的位置，而不能改变极点附近对 $p^2$ 的斜率。于是四维 $\varphi^4$ 理论在一圈阶满足
\begin{equation}
 \boxed{
 \delta m^2=O(g_4),
 \qquad
 Z_\varphi=1+O(g_4^2).}
 \label{eq:phi4-one-loop-mass-vs-Z-dp52}
\end{equation}
也就是说，质量反项在一圈就需要，场强反项要到二圈才第一次被二点函数要求。

四点函数的情况不同。$g_4^2$ 阶出现 $s,t,u$ 三个泡图，每个都依赖外部不变量，并具有同一个局域对数紫外极点。它们不能由质量反项消去，而必须重定义四次耦合。于是“哪一种反项出现”并不是预先背下来的规则，而是由 1PI 函数的外腿数和外部动量结构决定：二点函数的常数部分固定质量，二点函数的 $p^2$ 斜率固定场强，四点函数在指定运动学点的值固定耦合。

把重整化后二点函数写成
\begin{equation}
 \widetilde G_R(p)
 =\frac{\ii}
 {p^2-m_R^2c^2-\Sigma_R(p^2)+\ii0^+},
 \label{eq:phi4-ren-propagator-dp52}
\end{equation}
壳上方案直接把“物理质量”和“单位极点留数”翻译成
\begin{equation}
 \boxed{
 \Sigma_R(m_R^2c^2)=0,
 \qquad
 \left.\frac{\dd\Sigma_R}{\dd p^2}\right|_{p^2=m_R^2c^2}=0.}
 \label{eq:phi4-onshell-ren-conditions-dp52}
\end{equation}
第一式固定极点位置，第二式固定极点斜率。$\overline{\rm MS}$ 不直接采用这两个物理条件，而只减去维数正则化极点；两种方案得到的中间参数不同，但把参数都换成同一组实验输入后，可观测量只在被截掉的更高微扰阶上不同。


\medskip

\eqnrefs{eq:phi4-tadpole-selfenergy-dp52} 是 $\varphi^4$ 二点 1PI 函数的一圈蝌蚪贡献。由于积分与外部四动量无关，它在这一阶只移动传播子的极点位置，而不产生随 $p^2$ 变化的紫外斜率项。顶角收缩计数给出 $1/2$ 对称因子；维数正规化把紫外极点显式分离，质量与场强反项分别控制逆传播子的常数项和 $p^2$ 斜率。把重整化逆传播子在物理极点附近展开并固定极点位置与留数，便得到\eqnrefs{eq:phi4-onshell-ren-conditions-dp52}。



从相互作用
\begin{equation*}
 \Lag_{\rm int}
 =-\frac{g_4}{4!\,\hbar c}\varphi^4
\end{equation*}
的一阶 Dyson 插入开始。为形成二点 1PI 函数，两个外部场必须分别与顶角上的两个场收缩；选择有序外腿的方式数为
\begin{equation*}
 4\times3=12.
\end{equation*}
剩余两个顶角场彼此收缩只有一种方式。作用量中的 $1/4!=1/24$ 因而留下
\begin{equation*}
 \frac{12}{24}=\frac12,
\end{equation*}
这就是\eqnrefs{eq:phi4-tadpole-selfenergy-dp52} 的对称性因子。内部线从顶角出发又回到同一顶角，没有任何外部四动量流经圈，因此积分只依赖 $m_R$ 与正则化尺度，而不依赖 $p$。

令
\begin{equation*}
 T(m_R^2)
 \equiv
 \mu_{\rm DR}^{2\epsilon}
 \int\frac{\dd^d \ell}{(2\pi)^d}
 \frac{\ii}{\ell^2-m_R^2c^2+\ii0^+}.
\end{equation*}
用 Schwinger 参数
\begin{equation*}
 \frac{\ii}{\ell^2-m_R^2c^2+\ii0^+}
 =\int_0^\infty\dd\tau\,
 \ee^{\ii\tau(\ell^2-m_R^2c^2+\ii0^+)},
\end{equation*}
把动量积分化成 $d$ 维高斯积分，再对 $\tau$ 积分，得到解析延拓形式
\begin{equation*}
 T(m_R^2)
 =\frac{\mu_{\rm DR}^{2\epsilon}}
 {(4\pi)^{d/2}}
 \Gamma\!\left(1-\frac d2\right)
 (m_R^2c^2)^{d/2-1}.
\end{equation*}
取 $d=4-2\epsilon$，并展开
\begin{align*}
 \Gamma(-1+\epsilon)
 &=-\frac1\epsilon+\gamma_E-1+O(\epsilon),\\
 (4\pi)^\epsilon
 \left(\frac{\mu_{\rm DR}^2}{m_R^2c^2}\right)^\epsilon
 &=1+\epsilon
 \left[
 \ln4\pi+\ln\frac{\mu_{\rm DR}^2}{m_R^2c^2}
 \right]+O(\epsilon^2),
\end{align*}
所以
\begin{equation*}
 T(m_R^2)
 =-\frac{m_R^2c^2}{16\pi^2}
 \left[
 \frac1{\bar\epsilon}
 +1-\ln\frac{m_R^2c^2}{\mu_{\rm DR}^2}
 \right]
 +O(\epsilon),
\end{equation*}
其中 $1/\bar\epsilon=1/\epsilon-\gamma_E+\ln4\pi$。这个结果的关键结构是“局域、与 $p$ 无关、正比于 $m_R^2c^2$”；其有限常数部分会随减除方案改变，而这三个结构不会改变。

把裸二次部分写成重整化场与反项：
\begin{equation*}
 \frac12 Z_\varphi(\partial\varphi_R)^2
 -\frac12 Z_\varphi m_B^2c^2\frac{\varphi_R^2}{\hbar^2}
 =\frac12(\partial\varphi_R)^2
 -\frac12m_R^2c^2\frac{\varphi_R^2}{\hbar^2}
 +\Lag_{\rm ct}^{(2)}.
\end{equation*}
在动量空间，二点反项只能产生
\begin{equation*}
 \Sigma_{\rm ct}(p^2)
 =A_{\rm ct}+B_{\rm ct}(p^2-m_R^2c^2),
\end{equation*}
其中常数 $A_{\rm ct}$ 由质量反项控制，斜率 $B_{\rm ct}$ 由 $\delta Z_\varphi$ 控制。由于一圈蝌蚪满足
\begin{equation*}
 \frac{\dd\Sigma_{\rm tad}^{(1)}}{\dd p^2}=0,
\end{equation*}
一圈阶没有任何紫外斜率需要 $\delta Z_\varphi$ 抵消，因此
\begin{equation*}
 \delta Z_\varphi^{(1)}=0,
 \qquad
 Z_\varphi=1+O(g_4^2),
\end{equation*}
而常数紫外极点必须由 $\delta m^2=O(g_4)$ 吸收。这得到\eqnrefs{eq:phi4-one-loop-mass-vs-Z-dp52}。

最后检查壳上条件。精确重整化传播子定义为
\begin{equation*}
 \widetilde G_R(p)
 =\frac{\ii}
 {p^2-m_R^2c^2-\Sigma_R(p^2)+\ii0^+}.
\end{equation*}
要求物理极点仍位于 $p^2=m_R^2c^2$，分母在该点必须为零，因此
\begin{equation*}
 \Sigma_R(m_R^2c^2)=0.
\end{equation*}
在极点附近令 $p^2=m_R^2c^2+\delta$，Taylor 展开
\begin{align*}
 p^2-m_R^2c^2-\Sigma_R(p^2)
 &=\delta
 -\Sigma_R(m_R^2c^2)
 -\delta\,\Sigma_R'(m_R^2c^2)
 +O(\delta^2),\\
 &=\delta\,[1-\Sigma_R'(m_R^2c^2)]
 +O(\delta^2).
\end{align*}
若把重整化场定义成极点留数为 1，就必须再要求
\begin{equation*}
 \Sigma_R'(m_R^2c^2)=0.
\end{equation*}
这正是\eqnrefs{eq:phi4-onshell-ren-conditions-dp52}。因此质量反项与场强反项分别控制“极点位置”和“极点留数”；它们不是两个形式相似但物理意义不明的常数。
\end{exbox}

四维 \(\varphi^4\) 理论把“正则化—反项—尺度流”三步直接连在一起。一圈四点函数的三个通道都归结为同一个标量泡图。以 \(s\) 通道为例，先把需要计算的标量泡图明确写成
\begin{equation}
 \boxed{
 I(s)=\mu_{\rm DR}^{2\epsilon}
 \int\frac{\dd^d \ell}{(2\pi)^d}
 \frac{\ii}
 {\bigl(\ell^2-m_R^2c^2+\ii0^+\bigr)
  \bigl((\ell+p)^2-m_R^2c^2+\ii0^+\bigr)},
 \qquad p^2=s.}
 \label{eq:phi4-bubble-integral-r68}
\end{equation}
对\eqnrefs{eq:phi4-bubble-integral-r68} 做 Feynman 参数化和维数正则化后得到
\begin{equation}
 \boxed{
 I(s)
 =-\frac1{16\pi^2}
 \left[
 \frac1{\bar\epsilon}
 -\int_0^1\dd x\,
 \ln\frac{m_R^2c^2-x(1-x)s-\ii0^+}{\mu_{\rm DR}^2}
 \right]
 +\order(\epsilon),}
 \label{eq:phi4-bubble-result-dp10}
\end{equation}
其中 \(1/\bar\epsilon\equiv1/\epsilon-\gamma_E+\ln4\pi\)。三个通道具有相同的局域紫外极点，因此在 $d=4-2\epsilon$ 中，$\overline{\rm MS}$ 的一圈裸耦合可以写为
\begin{equation}
 \boxed{
 g_{4,B}
 =\mu^{2\epsilon}
 \left[
 g_4(\mu)
 +\frac{3}{32\pi^2}\frac{g_4^2(\mu)}{\bar\epsilon}
 +\order(g_4^3)
 \right].}
 \label{eq:phi4-bare-coupling-r68}
\end{equation}
裸参数不能依赖人为引入的减除尺度。对\eqnrefs{eq:phi4-bare-coupling-r68} 在固定 $g_{4,B}$ 下求尺度导数，得到
\begin{equation}
 \boxed{
 \beta_{g_4}(g_4)
 \equiv
 \mu\frac{\dd g_4}{\dd\mu}\Big|_{g_{4,B}}
 =\frac{3g_4^2}{16\pi^2}+\order(g_4^3).}
 \label{eq:phi4-beta-one-loop-dp10}
\end{equation}
一圈积分解为
\begin{equation}
 \boxed{
 \frac1{g_4(\mu)}
 =\frac1{g_4(\mu_0)}
 -\frac{3}{16\pi^2}\ln\frac\mu{\mu_0}.}
 \label{eq:phi4-running-one-loop-dp10}
\end{equation}

\begin{exbox}[一圈运行的数值尺度]
设某个参考尺度 $\mu_0$ 上测得 $g_4(\mu_0)=1$。先把分辨尺度提高三个数量级，考察一个仍然温和的尺度比
\begin{equation*}
 \frac\mu{\mu_0}=10^3,
 \qquad
 \ln\frac\mu{\mu_0}=6.9078.
\end{equation*}
一圈解给
\begin{align*}
 \frac1{g_4(\mu)}
 &=1-\frac{3}{16\pi^2}(6.9078)\\
 &\simeq0.8688,
\end{align*}
所以
\begin{equation*}
 g_4(10^3\mu_0)\simeq1.151.
\end{equation*}
耦合本身只增加约 $15\%$，但固定阶展开中真正控制大对数的是
\begin{equation*}
 \epsilon_{\log}
 \equiv\frac{g_4}{16\pi^2}\left|\ln\frac\mu{\mu_0}\right|
 \simeq4.37\times10^{-2},
\end{equation*}
仍明显小于 1，因此一圈结果在这个例子里仍是受控的。

继续把一圈解的分母外推到零，可得到形式上的 Landau 极点
\begin{equation*}
 \frac{\mu_L}{\mu_0}
 =\exp\!\left(\frac{16\pi^2}{3g_4(\mu_0)}\right)
 \simeq7.3\times10^{22}
 \qquad[g_4(\mu_0)=1].
\end{equation*}
这个巨大数字不应被误读成“一圈理论精确预言了一个真实新粒子能标”。它只说明正的一圈 $\beta$ 函数会让耦合随紫外尺度增大，并且若坚持把同一低阶方程外推足够远，最终一定离开微扰控制区。重整化群的价值正是把这种尺度积累显式量化出来。
\end{exbox}

普通固定阶一圈截断要求 \(g_4/(16\pi^2)\ll1\)；跨越很大的尺度比时还需 \(|g_4\ln(\mu/\mu_0)|/(16\pi^2)\ll1\)，否则应使用重整化群方程重求和大对数，而不是继续展开\eqnrefs{eq:phi4-bubble-result-dp10} 的有限阶对数。

任意重整化 $n$ 点函数的尺度依赖应从整个重整化参数空间一次建立。若裸场与重整化场满足 $\varphi_B=Z_\varphi^{1/2}\varphi_R$，同一组外点的裸与重整化关联函数满足
\begin{equation}
 \boxed{
 G_B^{(n)}=Z_\varphi^{n/2}G_R^{(n)}.}
 \label{eq:bare-ren-green-relation-r68}
\end{equation}
设理论含一组无量纲耦合 $\{g_a(\mu)\}$ 和一组质量参数 $\{m_b(\mu)\}$。耦合和质量对尺度的对数导数分别描述参数怎样随分辨尺度流动；场归一化因子 $Z_\varphi$ 的尺度导数则描述量子涨落使场偏离经典工程尺度律的程度。后一个量通常称为场的\emph{反常维数}；“反常”在这里只表示额外的量子尺度修正，并不表示理论发生了量子对称性反常。具体定义为
\begin{equation}
\boxed{
 \beta_{g_a}\equiv\mu\frac{\dd g_a}{\dd\mu}\bigg|_{\rm bare},
 \qquad
 \beta_{m_b}\equiv\mu\frac{\dd m_b}{\dd\mu}\bigg|_{\rm bare},
 \qquad
 \gamma_\varphi\equiv\frac12\mu\frac{\dd\ln Z_\varphi}{\dd\mu}\bigg|_{\rm bare}.}
\label{eq:rg-general-flow-definitions-dp74}
\end{equation}
固定所有裸参数时，$G_B^{(n)}$ 不依赖人为减除尺度，因此得到一般尺度方程
\begin{equation}
 \boxed{
 \left[
 \mu\frac{\partial}{\partial\mu}
 +\sum_a\beta_{g_a}\frac{\partial}{\partial g_a}
 +\sum_b\beta_{m_b}\frac{\partial}{\partial m_b}
 +n\gamma_\varphi
 \right]G_R^{(n)}=0.}
 \label{eq:callan-symanzik-general-dp10}
\end{equation}
这条方程描述重整化参数空间上的全尺度流。若只保留一个耦合 $g_{\rm R}$ 和一个质量 $m$，上式退化为
\[
 \left(
 \mu\partial_\mu+\beta_{g_{\rm R}}\partial_{g_{\rm R}}
 +\beta_m\partial_m+n\gamma_\varphi
 \right)G_R^{(n)}=0,
\]
这正是前面 $\varphi^4$ 一圈运行的参数特例。

局域复合算符具有另一种尺度结构。多个具有相同守恒量子数的算符在改变量子分辨尺度时可以彼此混合，因此它们的尺度修正不再由一个数描述，而由算符空间中的矩阵描述。这个矩阵就是复合算符的反常维数矩阵。若一组算符满足
\begin{equation}
\boxed{
 \mathcal O_i^B=Z_{ij}^{(\mathcal O)}\mathcal O_j^R,
 \qquad
 \gamma^{(\mathcal O)}
 =(Z^{(\mathcal O)})^{-1}\mu\frac{\dd Z^{(\mathcal O)}}{\dd\mu},}
\label{eq:operator-mixing-general-dp74}
\end{equation}
那么 $\gamma^{(\mathcal O)}$ 从一开始就是算符空间上的线性算符；其非对角元表示尺度流把一个基算符混入另一个基算符。只有在该空间恰为一维时，它才退化成单个反常维数。

\phantomsection\label{sec:qft-eft-unitarity-dp10}
重整化控制给定理论内部的紫外敏感性，却不要求低能实验显式解析所有高能自由度。若某些自由度的特征能量远高于实验能量，它们不能作为真实外态产生，但仍会通过虚过程改变低能振幅。

有效场论（effective field theory，EFT）只保留当前分辨率能够直接激发的轻自由度，并把重自由度的虚效应编码为由轻场构成的局域高维算符。若实验特征动量远小于重尺度，局域展开就按“实验动量/重尺度”的幂次有序；算符维数越高，通常受到越高次幂的抑制。EFT 的适用范围因此首先由这一无量纲展开参数控制，并还需满足散射矩阵幺正性等一致性界限。

采用 \(x^0=ct\)，定义几何拉氏量密度 \(\widehat\Lag=\Lag/(\hbar c)\)。同一个重尺度可等价写成逆长度、物理动量或能量：
\begin{equation}
 \boxed{
 \Lambda_*\ [\mathrm{m^{-1}}],
 \qquad
 p_*\equiv\hbar\Lambda_*,
 \qquad
 \Lambda_E\equiv cp_*=\hbar c\Lambda_*.}
 \label{eq:qft-eft-three-scales-dp42}
\end{equation}
若局域算符 \(\widehat{\mathcal O}_i\) 的逆长度工程维数为 \(\Delta_i>4\)，一般低能有效拉氏量为
\begin{equation}
 \boxed{
 \widehat\Lag_{\rm EFT}
 =\widehat\Lag_{\Delta\le4}
 +\sum_{\Delta_i>4}
 \frac{c_i}{\Lambda_*^{\Delta_i-4}}
 \widehat{\mathcal O}_i,}
 \label{eq:eft-operator-suppression-dp10}
\end{equation}
其中 \(c_i\) 无量纲。它们记录被消去高能自由度的短距离信息，而低能算符 \(\widehat{\mathcal O}_i\) 描述仍可分辨的自由度；乘在局域有效算符前的这些短距离系数称为\emph{Wilson 系数}。

Wilson 系数不能仅凭低能理论自身决定：应选择一组低能外态或 Green 函数，在包含重自由度的完整理论与有效理论中分别计算，并要求两种描述在共同适用的低能区域逐阶给出相同结果。用这种“比较同一低能物理并据此固定 Wilson 系数”的步骤称为\emph{匹配}（matching）。匹配解决的是高能信息怎样进入低能系数，而随后随重整化尺度改变这些系数则是另一个独立的重整化群问题。特征外部物理动量为 \(Q\) 时，维数 \(\Delta_i\) 的算符典型贡献按
\begin{equation}
 \boxed{
 \epsilon_{\rm EFT}^{(i)}
 \sim
 \left(\frac{Q}{p_*}\right)^{\Delta_i-4}
 \ll1}
 \label{eq:eft-power-counting-dp10}
\end{equation}
排序。若 \(Q\sim p_*\)，局域高维展开失去层级，必须恢复被消去的高能自由度或改用非局域描述。

有效场论把未解析的高能自由度编码为高维局域算符，并以 \(Q/p_*\) 排序；Feshbach--Schur、Schrieffer--Wolff 与 Green 张量投影则利用谱隙或预解式消去远失谐或连续自由度，控制参数通常是耦合与能隙之比以及预解式的能量依赖。二者都只保留给定能窗内可分辨的自由度，但展开参数不同。


有效场论截断还必须满足散射矩阵幺正性。对无自旋、单通道弹性 \(2\to2\) 散射，在质心系定义
\begin{equation}
 \boxed{
 \mathcal M(s,\theta)
 =16\pi\sum_{\ell=0}^{\infty}(2\ell+1)
 a_\ell(s)P_\ell(\cos\theta).}
 \label{eq:partial-wave-expansion-dp10}
\end{equation}
单弹性通道中 \(S_\ell=1+2\ii a_\ell\)，由 \(|S_\ell|=1\) 得到 Argand 圆，因此
\begin{equation}
 \boxed{
 |a_\ell|\le1,
 \qquad
 |\Re a_\ell|\le\frac12.}
 \label{eq:partial-wave-unitarity-bound-dp10}
\end{equation}
存在非弹性通道时允许区域收缩到圆盘内部。若树级有效场论振幅已经增长到 \(|\Re a_\ell|\gtrsim1/2\)，失效的是固定阶低能截断；此时必须加入圈修正、重求和或新的高能自由度，而不能把精确理论解释为“不幺正”。

幂次计数判断潜在紫外行为；正则化定义发散积分；重整化用有限输入定义参数；重整化群组织减除尺度变化；有效场论按 \(Q/p_*\) 排列未解析高能结构；分波幺正性给出振幅的非微扰边界。QED、QCD 与电弱理论的规范群和场内容不同，但共享这套紫外组织原则。


这些结构可与 \cite{PeskinSchroeder1995,Schwartz2014QFT} 对照；全部尺度均按物理四动量与 SI 约定解释，并显式保留 \(c\) 与 \(\hbar\)。
